We conduct a series of experiments across multiple tasks, including reading comprehension, mathematics, reasoning, and multitasking. These experiments are designed to evaluate \framework{}’s self-improvement capabilities in comparison to both hand-designed agents and a state-of-the-art automated agent design method.
In addition, to gain deeper insights into the behavior and performance of \framework, we also conduct a case study with Game of 24 as presented in Section \ref{sec:game24}.

\footnotetext{The results of baseline models are refer to \citet{hu2024automated}.}

\subsection{Baseline Methods}
To establish a comprehensive baseline, we select both hand-designed methods and automated agent design techniques. Hand-designed methods are well-known approaches that include:
1) Chain-of-Thought (CoT) \citep{wei2022chain} that encourages agents to reason step-by-step before providing an answer.
2) Self-Consistency with CoT (CoT-SC) \citep{wang2023selfconsistencyimproveschainthought} that generates multiple solution paths using CoT and selects the most consistent answer.
3) Self-Refine \citep{madaan2024self} that involves agents assessing their outputs and correcting mistakes in subsequent attempts.
4) LLM-Debate \citep{du2023improvingfactualityreasoninglanguage} that allows different LLMs to engage in a debate, offering diverse viewpoints.
5) Step-back Abstraction \citep{zheng2024stepbackevokingreasoning} that prompts agents to initially focus on fundamental principles before diving into task details.
6) Quality-Diversity \citep{lu2024aiscientistfullyautomated} that generates diverse solutions and combines them.
7) Role Assignment \citep{xu2023expertpromptinginstructinglargelanguage} that assigns specific roles to LLMs to generate better solutions by leveraging different perspectives.
Given the limitations of fixed algorithms in handling dynamic scenarios, we select 8) Meta Agent Search \citep{hu2024automated}, the latest state-of-the-art method for automated agent design, as our main comparison point. 
%Meta Agent Search represents advanced approach for dynamic agent improvement and is therefore a suitable baseline for evaluating \framework.

\subsection{Experimental Settings}
Following the setup of \citet{hu2024automated}, we evaluate \framework's self-improvement capabilities across four well-known benchmarks:
1) DROP \citep{dua2019dropreadingcomprehensionbenchmark} for reading comprehension.
2) MGSM \citep{shi2022languagemodelsmultilingualchainofthought} for testing mathematical skills in a multilingual context.
3) MMLU \citep{hendrycks2021measuringmassivemultitasklanguage} for evaluating multi-task problem-solving abilities.
4) GPQA \citep{rein2023gpqagraduatelevelgoogleproofqa} for tackling challenging graduate-level science questions.

% Given the complexity of the tasks and the need for advanced reasoning and understanding, the improvement cycle of \framework is driven by GPT-4o. 
% In the main experiment, we implement two different settings: 1) To make a fair comparison with baseline methods, we forbid \framework to change the API of the LLM used to perform the tasks (by default GPT-3.5) and use a closed-book approach with no access to the Internet, and 2) To explore the upper bound of \framework's capabilities, we remove all constraints. 
Given its simplicity and versatility, we use CoT as the initial policy for all tasks.
In addition, as shown in Section \ref{sec:game24}, we also analyze the performance of \framework when using other algorithms as the initial policies.




We perform 6 independent self-improvement cycles on the validation dataset for each task, with a maximum of 30 iterations per cycle. Each cycle represents a complete self-improvement process, where \framework iteratively modifies its logic to enhance performance.
After obtaining the optimized agent, we test it on the test set. For fairness, we use GPT-3.5 for all the tests, whether for the baseline or \framework{}.
Further details can be found in Appendix \ref{sec:expdetails}.
%regarding the experimental setup and additional results .


\subsection{Experimental Results and Analysis}
The experimental results are shown in Table \ref{tab:main}. Under the same setting, \framework achieves either optimal or comparable results to Meta Agent Search across all tasks. Notably, in the mathematics task MGSM, \framework outperforms it by 11\%. This suggests that reasoning tasks offer greater room for improvement for \framework (\textit{performance}). 
In contrast to Meta Agent Search, which needs to design different modules for different tasks, \framework demonstrates greater \textit{flexibility}. It requires only a simple initial policy, such as CoT, with all other components being autonomously generated. 
Moreover, through interaction with the environment, it gradually adapts and independently devises effective methods for the current task. The final policies generated by \framework are shown in Appendix \ref{sec:maincode}. Additionally, our method converges faster, with the required number of iterations and computational \textit{cost} compared to the Meta Agent shown in Appendix \ref{sec:cost}. 
% In addition, the optimization curve of \framework is shown in Figure \ref{}. We can see that it converges rapidly and demonstrates the ability to quickly self-correct when errors occur.

We also conduct experiments without restrictions, where \framework significantly outperforms all baselines. Upon further analysis, we find that this is primarily due to the agent's spontaneous requests for assistance from more powerful models such as GPT-4o in some tasks. Therefore, \framework is particularly well-suited for open-ended scenarios, where it can employ various strategies to enhance performance (\textit{potential}).

Therefore, we can find that \framework{} is superior to the previous agent frameworks in terms of performance, flexibility, cost, and potential.